% ====================== TÓM TẮT ======================
\chapter*{Tóm tắt}
\addcontentsline{toc}{chapter}{Tóm tắt}

Phân đoạn ngữ nghĩa (\textit{semantic segmentation}) ảnh chụp từ thiết bị bay
không người lái (UAV) là bài toán nền tảng cho các ứng dụng giám sát và kiểm kê
rừng. Tuy nhiên, dữ liệu rừng thu thập từ UAV có đặc trưng \textbf{mất cân bằng
lớp nghiêm trọng}: trên bộ dữ liệu sử dụng trong báo cáo, tỷ lệ giữa lớp phổ biến
nhất và lớp hiếm nhất lên tới \textbf{2.761 lần} (Ground\_vegetation $32{,}4\%$ so
với Fence $0{,}01\%$). Hệ quả là các mô hình huấn luyện bằng hàm mất mát
Cross-Entropy thuần có xu hướng bỏ qua các lớp hiếm nhưng quan trọng như
\textit{Fence}, \textit{Car}, \textit{Fallen\_trees}.

Báo cáo này nghiên cứu một cách hệ thống bốn chiến lược xử lý mất cân bằng lớp
trên kiến trúc \textbf{UNet++ với encoder ResNet50}: (1) Cross-Entropy làm chuẩn
đối sánh; (2) Cross-Entropy có trọng số lớp; (3) hàm mất mát kết hợp
Focal\,+\,Dice\,+\,Cross-Entropy với trọng số cố định; và (4) cùng tổ hợp đó nhưng
trọng số được \textit{học thích nghi} theo cơ chế bất định đa nhiệm (uncertainty
weighting). Kết quả thực nghiệm cho thấy chiến lược thích nghi đạt
\textbf{mIoU $= 0{,}8365$}, cao hơn chuẩn Cross-Entropy ($0{,}8187$) $1{,}78$ điểm
phần trăm, và quan trọng hơn, cải thiện mạnh trên các lớp hiếm
(Fence: $+10{,}9$ điểm; Car: $+4{,}4$ điểm), trong khi vẫn giữ được hiệu năng trên
các lớp phổ biến.

Báo cáo cũng trình bày cơ sở lý thuyết của từng hàm mất mát, quy trình thực nghiệm,
phân tích lỗi theo lớp, và các hướng mở rộng đang được tiến hành (đánh giá trên các
điều kiện thời tiết \textit{sunny}, \textit{overcast}, \textit{sunny+overcast} và bộ
dữ liệu \textit{WildUAV}).

\vspace{0.5cm}
\noindent\textbf{Từ khóa:} phân đoạn ngữ nghĩa, ảnh UAV, giám sát rừng, mất cân bằng
lớp, Focal Loss, Dice Loss, UNet++, học sâu.

% ----- (Tùy chọn) Abstract tiếng Anh -----
\TBD{Nếu trường yêu cầu, bổ sung bản Abstract tiếng Anh tại đây.}
